\chapter{Higher-Order Deforestation and Defunctionalization}
\label{chap:higher_order}

\section{The Challenge of Higher-Order Driving}
\label{sec:higher_order:challenge}

In functional languages, higher-order functions (closures that capture runtime environments) are the primary vehicle of abstraction. However, specializing higher-order code is notoriously difficult:
\begin{enumerate}
    \item Closures hide control flow behind indirect function pointers.
    \item Closure environments capture heap-allocated records, preventing syntactic term matching.
    \item Mutual recursion through closure arguments easily induces whistle blowouts.
\end{enumerate}

\section{Whole-Program Reynolds Defunctionalization (Phase 49)}
\label{sec:higher_order:reynolds}

NumLang conquers this challenge by implementing a whole-program \emph{Reynolds Defunctionalization} pass in \texttt{src/mir/defunctionalize.rs} \citep{reynolds1972definitional}.

\subsection{Transformation Pipeline}
\begin{enumerate}
    \item \textbf{Call Signature Analysis}: The compiler collects all lambda expressions and function pointer types across the program.
    \item \textbf{ClosureTag Synthesis}: For each unique function signature $\tau_{\text{sig}} = \mathtt{fn}(\tau_1, \dots, \tau_n) \to \tau_{\text{ret}}$, NumLang generates a strongly typed enum:
    \begin{lstlisting}[language=NumLang]
    enum ClosureTag_Sig {
        Lambda_1(i64),       // captures single i64
        Lambda_2(Box<List>), // captures heap list
    }
    \end{lstlisting}
    \item \textbf{Lowering Allocations}: \texttt{ClosureAlloc} statements are replaced by enum constructor initializations.
    \item \textbf{Dispatch Switch Synthesis}: All indirect calls $\mathtt{IndirectCall}(fn\_ptr, \vec{args})$ are replaced by explicit \texttt{Terminator::Switch} dispatchers branching on the closure discriminant tag.
    \item \textbf{Scalar Replacement of Aggregates (SROA)}: Environment payloads are unpacked into primitive SSA scalar registers.
\end{enumerate}

Following defunctionalization, the entire program is strictly first-order SSA MIR.

\section{Symbolic Driving Through Closures (Phase 34)}
\label{sec:higher_order:driving_closures}

Because defunctionalization exposes closure tags as regular enums, the supercompiler's existing pattern matching driving mechanisms handle closures seamlessly:
\begin{itemize}
    \item Symbolic closures are tracked as $\mathtt{SymTerm::ClosureVal}(f, \vec{captured})$.
    \item When a call site invokes the closure, driving resolves the target statically and inlines the function body with the captured arguments bound in the local environment.
    \item Dynamic dispatch overhead is eliminated completely.
\end{itemize}

\section{Lazy Codata Driving and Stream Fusion (Phase 51)}
\label{sec:higher_order:codata}

To achieve true GHC-style stream fusion without allocating intermediate iterator structures, NumLang implements \emph{Lazy Codata Driving} (\texttt{src/mir/thunk\_analysis.rs}).

Pipelines over infinite streams (such as \texttt{iterate}, \texttt{map}, and \texttt{take}) are expressed via suspended thunks:
\begin{equation}
x = \mathtt{Rvalue::Thunk}(f, \vec{args}), \quad y = \mathtt{Terminator::Force}(x)
\end{equation}
The supercompiler propagates demand backwards: a thunk is only driven when forced. When a pipeline consists of chained thunk forces, driving eliminates the suspension overhead, fusing the generator, transformer, and consumer into a single flat loop with zero heap allocations.

\section{The Reynolds Defunctionalization Engine}
\label{sec:higher_order:defunc_impl}

Listing~\ref{lst:defunc_impl} details the Reynolds defunctionalization pass from \texttt{src/mir/defunctionalize.rs}, replacing higher-order closures with first-order sum types:

\begin{lstlisting}[language=Rust, caption={Reynolds Defunctionalization Algorithm (\texttt{src/mir/defunctionalize.rs}).}, label={lst:defunc_impl}]
//!
//! Compiles higher-order closures and indirect calls into zero-allocation,
//! monomorphic, first-order SSA forms with static dispatch:
//!
//! 1. Discovers all closures (`Rvalue::ClosureAlloc`) and groups them by call signature.
//! 2. Synthesizes a discriminated union enum `ClosureTag_<Signature>` per distinct signature.
//! 3. Rewrites `Rvalue::ClosureAlloc` into typed `Rvalue::EnumVariant` carrying captured environment payloads.
//! 4. Lowers `Terminator::IndirectCall` into direct `Terminator::Switch` over tags, dispatching
//!    to monomorphic static `Rvalue::Call` sites in specialized basic blocks.

use std::collections::BTreeMap;
use crate::span::Span;
use crate::mir::{BasicBlockId, Place, Projection, Terminator};
use crate::mir::lower::{MirBasicBlock, MirLocalDecl, MirProgram, Rvalue, Statement};
use crate::typecheck::typed_ast::{TypedEnumDef, TypedEnumVariant};
use crate::typecheck::types::Type;

#[derive(Debug, Clone, PartialEq, Eq, Hash)]
pub struct ClosureSignature {
    pub param_tys: Vec<Type>,
    pub return_ty: Type,
}

#[derive(Debug, Clone)]
struct ClosureCandidate {
    fn_name: String,
    captured_types: Vec<Type>,
}

#[derive(Debug, Clone)]
struct ClosureGroup {
    signature: ClosureSignature,
    candidates: Vec<ClosureCandidate>,
}

#[derive(Debug, Clone)]
struct ResolvedVariant {
    fn_name: String,
    tag: usize,
    captured_types: Vec<Type>,
}

#[derive(Debug, Clone)]
struct SignatureGroup {
    signature: ClosureSignature,
    variants: Vec<ResolvedVariant>,
}

#[derive(Debug, Clone, Default, PartialEq, Eq)]
pub struct DefunctionalizeStats {
    pub closures_defunctionalized: usize,
    pub indirect_calls_resolved: usize,
    pub synthetic_enums_created: usize,
}

/// Executes whole-program Reynolds defunctionalization on the given MIR program.
pub fn defunctionalize_program(program: &mut MirProgram) -> DefunctionalizeStats {
    let mut stats = DefunctionalizeStats::default();

    // 1. Scan for all closure allocations to discover signatures and variants
    let mut closure_allocs: Vec<(String, usize)> = Vec::new(); // (fn_name, captured_count)
    for func in &program.functions {
        for block in &func.blocks {
            for stmt in &block.statements {
                let Statement::Assign(_, rval) = stmt;
                if let Rvalue::ClosureAlloc { fn_name, captured } = rval {
                    if !closure_allocs.iter().any(|(n, _)| n == fn_name) {
                        closure_allocs.push((fn_name.clone(), captured.len()));
                    }
                }
            }
        }
    }

    if closure_allocs.is_empty() {
        return stats;
    }

    // 2. Map closure functions to their callable signatures and captured types
    let mut groups: Vec<ClosureGroup> = Vec::new();
    for (fn_name, cap_count) in &closure_allocs {
        if let Some(target_fn) = program.functions.iter().find(|f| &f.name == fn_name) {
            let cap_tys: Vec<Type> = target_fn.params[0..*cap_count]
                .iter()
                .map(|(_, ty)| ty.clone())
                .collect();
            let callable_tys: Vec<Type> = target_fn.params[*cap_count..]
                .iter()
                .map(|(_, ty)| ty.clone())
                .collect();
            let sig = ClosureSignature {
                param_tys: callable_tys,
                return_ty: target_fn.return_ty.clone(),
            };
            let cand = ClosureCandidate {
                fn_name: fn_name.clone(),
                captured_types: cap_tys,
            };
            if let Some(grp) = groups.iter_mut().find(|g| g.signature == sig) {
                grp.candidates.push(cand);
            } else {
                groups.push(ClosureGroup {
                    signature: sig,
                    candidates: vec![cand],
                });
            }
        }
    }

    // 3. Synthesize discriminated union sum types (enums)
    // Map fn_name -> (enum_name, variant_name, tag, captured_tys)
    let mut variant_map: BTreeMap<String, (String, String, usize, Vec<Type>)> = BTreeMap::new();
    let mut sig_groups: Vec<SignatureGroup> = Vec::new();

    for (sig_idx, group) in groups.iter().enumerate() {
        let enum_name = format!("ClosureTag_{}_{}", sig_idx, group.signature.param_tys.len());

        let mut enum_variants = Vec::new();
        let mut resolved_variants = Vec::new();
        for (tag, cand) in group.candidates.iter().enumerate() {
            let variant_name = format!("Variant_{}", cand.fn_name);
            enum_variants.push(TypedEnumVariant {
                name: variant_name.clone(),
                tag,
                payload: cand.captured_types.clone(),
                span: Span { start: 0, end: 0 },
            });
            variant_map.insert(
                cand.fn_name.clone(),
                (enum_name.clone(), variant_name, tag, cand.captured_types.clone()),
            );
            resolved_variants.push(ResolvedVariant {
                fn_name: cand.fn_name.clone(),
                tag,
                captured_types: cand.captured_types.clone(),
            });
            stats.closures_defunctionalized += 1;
        }

        sig_groups.push(SignatureGroup {
            signature: group.signature.clone(),
            variants: resolved_variants,
        });

        program.enums.push(TypedEnumDef {
            name: enum_name,
            variants: enum_variants,
            span: Span { start: 0, end: 0 },
        });
        stats.synthetic_enums_created += 1;
    }

    // 4. Rewrite Rvalue::ClosureAlloc into Rvalue::EnumVariant
    for func in &mut program.functions {
        for block in &mut func.blocks {
            for stmt in &mut block.statements {
                let Statement::Assign(_, rval) = stmt;
                if let Rvalue::ClosureAlloc { fn_name, captured } = rval {
                    if let Some((enum_name, variant_name, tag, _)) = variant_map.get(fn_name) {
                        *rval = Rvalue::EnumVariant {
                            enum_name: enum_name.clone(),
                            variant_name: variant_name.clone(),
                            tag: *tag,
                            fields: captured.clone(),
                        };
                    }
                }
            }
        }
    }

    // 5. Rewrite Terminator::IndirectCall into direct static dispatch Switch
    for func in &mut program.functions {
        let mut new_blocks: Vec<MirBasicBlock> = Vec::new();
        let mut max_bb = func.blocks.iter().map(|b| b.id.0).max().unwrap_or(0) + 1;

        for block in &mut func.blocks {
            if let Terminator::IndirectCall { callee, args, dest, next } = &block.terminator {
                // Find matching signature by arity
                let matching_group = sig_groups.iter().find(|g| {
                    g.signature.param_tys.len() == args.len()
                });

                if let Some(group) = matching_group {
                    let tag_place = Place {
                        local: format!("_dt_tag_{}_{}", callee.local, block.id.0),
                        projections: vec![],
                    };
                    func.locals.push(MirLocalDecl {
                        name: tag_place.local.clone(),
                        ty: Type::I64,
                        mutable: true,
                    });

                    // Tag read statement via Discriminant
                    block.statements.push(Statement::Assign(
                        tag_place.clone(),
                        Rvalue::Discriminant(callee.clone()),
                    ));

                    let mut targets = Vec::new();
                    for variant in &group.variants {
                        let arm_bb_id = BasicBlockId(max_bb);
                        max_bb += 1;
                        let mut arm_stmts = Vec::new();
                        let mut call_args = Vec::with_capacity(variant.captured_types.len() + args.len());

                        // Unpack captured environment payload fields
                        for (j, cap_ty) in variant.captured_types.iter().enumerate() {
                            let cap_temp = Place {
                                local: format!("_dt_cap_{}_{}_{}", variant.fn_name, j, arm_bb_id.0),
                                projections: vec![],
                            };
                            func.locals.push(MirLocalDecl {
                                name: cap_temp.local.clone(),
                                ty: cap_ty.clone(),
                                mutable: true,
                            });
                            let mut proj_place = callee.clone();
                            proj_place.projections.push(Projection::Payload(j));
                            arm_stmts.push(Statement::Assign(cap_temp.clone(), Rvalue::Use(proj_place)));
                            call_args.push(cap_temp);
                        }

                        // Pass remaining args
                        for arg in args {
                            call_args.push(arg.clone());
                        }

                        // Direct monomorphic call
                        arm_stmts.push(Statement::Assign(
                            dest.clone(),
                            Rvalue::Call(variant.fn_name.clone(), call_args),
                        ));

                        new_blocks.push(MirBasicBlock {
                            id: arm_bb_id.clone(),
                            arguments: vec![],
                            statements: arm_stmts,
                            terminator: Terminator::Branch { target: next.clone() },
                        });

                        targets.push((variant.tag as i64, arm_bb_id));
                    }

                    let default_bb_id = BasicBlockId(max_bb);
                    max_bb += 1;
                    new_blocks.push(MirBasicBlock {
                        id: default_bb_id.clone(),
                        arguments: vec![],

\end{lstlisting}

\section{Stream Fusion and Codata Thunk Analysis}
\label{sec:higher_order:thunk_impl}

Listing~\ref{lst:thunk_impl} reproduces the thunk analysis pass from \texttt{src/mir/thunk\_analysis.rs}:

\begin{lstlisting}[language=Rust, caption={Codata Thunk Analysis and Stream Fusion (\texttt{src/mir/thunk\_analysis.rs}).}, label={lst:thunk_impl}]
//!
//! Computes which thunks are demanded on execution paths through a MIR function:
//! - `Bottom`: thunk is never forced on any reachable path.
//! - `GuardDemanded`: thunk is forced on some paths (e.g. inside a conditional or switch branch).
//! - `FullyDemanded`: thunk is unconditionally forced on every execution path to exit.

use std::collections::{HashMap, HashSet};
use crate::mir::{BasicBlockId, Terminator};
use crate::mir::lower::{MirFunction, Rvalue, Statement};

#[derive(Debug, Clone, Copy, PartialEq, Eq, PartialOrd, Ord, Hash)]
pub enum Demand {
    /// Thunk is never forced on any path through this function.
    Bottom,
    /// Thunk is forced only on some paths (e.g. inside a conditional).
    GuardDemanded,
    /// Thunk is forced on every execution path (unconditionally demanded).
    FullyDemanded,
}

impl Demand {
    pub fn join(self, other: Demand) -> Demand {
        match (self, other) {
            (Demand::Bottom, d) | (d, Demand::Bottom) => d,
            (Demand::FullyDemanded, Demand::FullyDemanded) => Demand::FullyDemanded,
            _ => Demand::GuardDemanded,
        }
    }

    pub fn branch_meet(self, other: Demand) -> Demand {
        match (self, other) {
            (Demand::FullyDemanded, Demand::FullyDemanded) => Demand::FullyDemanded,
            (Demand::Bottom, Demand::Bottom) => Demand::Bottom,
            _ => Demand::GuardDemanded,
        }
    }
}

/// Computes thunk demand classification for all thunks in a MIR function.
pub fn compute_thunk_demands(func: &MirFunction) -> HashMap<String, Demand> {
    if func.blocks.is_empty() {
        return HashMap::new();
    }

    let mut candidate_thunks: HashSet<String> = HashSet::new();
    for block in &func.blocks {
        for stmt in &block.statements {
            let Statement::Assign(dest, rval) = stmt;
            if let Rvalue::Thunk { .. } = rval {
                candidate_thunks.insert(dest.local.clone());
            }
        }
        if let Terminator::Force { thunk, .. } = &block.terminator {
            candidate_thunks.insert(thunk.clone());
        }
    }

    if candidate_thunks.is_empty() {
        return HashMap::new();
    }

    let mut in_demands: HashMap<BasicBlockId, HashMap<String, Demand>> = HashMap::new();
    let mut out_demands: HashMap<BasicBlockId, HashMap<String, Demand>> = HashMap::new();

    for block in &func.blocks {
        let mut initial = HashMap::new();
        for t in &candidate_thunks {
            initial.insert(t.clone(), Demand::Bottom);
        }
        in_demands.insert(block.id.clone(), initial.clone());
        out_demands.insert(block.id.clone(), initial);
    }

    let mut changed = true;
    let mut iterations = 0;
    let max_iterations = func.blocks.len() * candidate_thunks.len() * 3 + 16;

    while changed && iterations < max_iterations {
        changed = false;
        iterations += 1;

        for block in func.blocks.iter().rev() {
            // Compute OUT[b] from successors' IN
            let mut new_out: HashMap<String, Demand> = HashMap::new();
            for t in &candidate_thunks {
                new_out.insert(t.clone(), Demand::Bottom);
            }

            match &block.terminator {
                Terminator::Return { .. } | Terminator::Unreachable => {
                    // Stays Bottom
                }
                Terminator::Branch { target } => {
                    if let Some(target_in) = in_demands.get(target) {
                        for (t, d) in target_in {
                            new_out.insert(t.clone(), *d);
                        }
                    }
                }
                Terminator::Force { cont, .. } => {
                    if let Some(cont_in) = in_demands.get(cont) {
                        for (t, d) in cont_in {
                            new_out.insert(t.clone(), *d);
                        }
                    }
                }
                Terminator::BranchIf { then_target, else_target, .. } => {
                    let default_in = HashMap::new();
                    let then_in = in_demands.get(then_target).unwrap_or(&default_in);
                    let else_in = in_demands.get(else_target).unwrap_or(&default_in);

                    for t in &candidate_thunks {
                        let d_then = then_in.get(t).copied().unwrap_or(Demand::Bottom);
                        let d_else = else_in.get(t).copied().unwrap_or(Demand::Bottom);
                        new_out.insert(t.clone(), d_then.branch_meet(d_else));
                    }
                }
                Terminator::Switch { targets, default, .. } => {
                    let default_in = HashMap::new();
                    let def_in = in_demands.get(default).unwrap_or(&default_in);

                    for t in &candidate_thunks {
                        let mut all_demands = Vec::with_capacity(targets.len() + 1);
                        all_demands.push(def_in.get(t).copied().unwrap_or(Demand::Bottom));
                        for (_, target_bb) in targets {
                            let t_in = in_demands.get(target_bb).unwrap_or(&default_in);
                            all_demands.push(t_in.get(t).copied().unwrap_or(Demand::Bottom));
                        }

                        let combined = if all_demands.iter().all(|&d| d == Demand::FullyDemanded) {
                            Demand::FullyDemanded
                        } else if all_demands.iter().any(|&d| d != Demand::Bottom) {
                            Demand::GuardDemanded
                        } else {
                            Demand::Bottom
                        };
                        new_out.insert(t.clone(), combined);
                    }
                }
                Terminator::IndirectCall { next, .. } => {
                    if let Some(next_in) = in_demands.get(next) {
                        for (t, d) in next_in {
                            new_out.insert(t.clone(), *d);
                        }
                    }
                }
                Terminator::Fork { left, right, .. } => {
                    let default_in = HashMap::new();
                    let l_in = in_demands.get(left).unwrap_or(&default_in);
                    let r_in = in_demands.get(right).unwrap_or(&default_in);
                    for t in &candidate_thunks {
                        let dl = l_in.get(t).copied().unwrap_or(Demand::Bottom);
                        let dr = r_in.get(t).copied().unwrap_or(Demand::Bottom);
                        new_out.insert(t.clone(), dl.join(dr));
                    }
                }
                Terminator::TypeGuard { fast_path, deopt_stub, .. } => {
                    let default_in = HashMap::new();
                    let fast_in = in_demands.get(fast_path).unwrap_or(&default_in);
                    let deopt_in = in_demands.get(deopt_stub).unwrap_or(&default_in);
                    for t in &candidate_thunks {
                        let d_fast = fast_in.get(t).copied().unwrap_or(Demand::Bottom);
                        let d_deopt = deopt_in.get(t).copied().unwrap_or(Demand::Bottom);
                        new_out.insert(t.clone(), d_fast.branch_meet(d_deopt));
                    }
                }
            }

            // Transfer function: IN[b] = OUT[b] updated by block body and terminator
            let mut new_in = new_out.clone();
            if let Terminator::Force { thunk, .. } = &block.terminator {
                new_in.insert(thunk.clone(), Demand::FullyDemanded);
            }

            let curr_in = in_demands.get_mut(&block.id).expect("block exists");
            if *curr_in != new_in {
                *curr_in = new_in;
                changed = true;
            }
            out_demands.insert(block.id.clone(), new_out);

\end{lstlisting}
